% TOP_PROBLEM: {"id": 1097, "title": "N-Queens Problem Asymptotics", "category_id": 2, "category": {"id": 2, "name": "combinatorics", "display_name": "Combinatorics", "description": "Counting problems, graph theory, discrete structures.", "slug": "combinatorics", "order_index": 2, "created_at": "2026-07-31T15:26:25.671Z"}, "status": "open"}
% FIRST_POSTED: null
% ENTRY_KIND: statement_only
% Imported from: unsolvedmath_additional_500.tex; UnsolvedMath ID 1097
% !TeX program = lualatex
% Compile twice: lualatex 1097.tex
% Self-contained: no external bibliography, images, or \input files.
% Numeric section labels and visible numbers are original UnsolvedMath IDs.
\documentclass[11pt,a4paper]{article}
\usepackage[margin=24mm,headheight=14pt]{geometry}
\usepackage{fontspec}
\setmainfont{Latin Modern Roman}
\setsansfont{Latin Modern Sans}
\setmonofont{Latin Modern Mono}
\usepackage{amsmath,amssymb,mathtools}
\usepackage{unicode-math}
\setmathfont{Latin Modern Math}
\usepackage{microtype}
\usepackage{enumitem,xurl,needspace}
\usepackage[dvipsnames]{xcolor}
\usepackage{hyperref}
\hypersetup{unicode=true,colorlinks=true,linkcolor=MidnightBlue,urlcolor=MidnightBlue,
 pdftitle={UnsolvedMath Problem 1097},pdfauthor={Source-annotated compilation},bookmarksnumbered=true}
\usepackage{bookmark,tocloft,titlesec,fancyhdr}
\setlength{\cftsecnumwidth}{4.2em}
\cftsetpnumwidth{2.5em}
\cftsetrmarg{4em}
\titleformat{\section}{\Large\bfseries\raggedright}{\thesection}{0.7em}{}
\pagestyle{fancy}\fancyhf{}
\fancyhead[L]{\small\sffamily Additional UnsolvedMath statements}
\fancyhead[R]{\small Original catalogue IDs}
\fancyfoot[C]{\thepage}
\setlength{\parindent}{0pt}
\setlength{\parskip}{5pt plus 1pt}
\setlength{\emergencystretch}{3em}
\setcounter{tocdepth}{1}\setcounter{secnumdepth}{1}
\setlist[itemize]{leftmargin=3em,itemsep=4pt,topsep=4pt,parsep=0pt}
\allowdisplaybreaks
\newcommand{\N}{\mathbb{N}}
\newcommand{\Z}{\mathbb{Z}}
\newcommand{\Q}{\mathbb{Q}}
\newcommand{\R}{\mathbb{R}}
\newcommand{\C}{\mathbb{C}}
\newcommand{\F}{\mathbb{F}}
\newcommand{\A}{\mathbb{A}}
\newcommand{\PP}{\mathbb{P}}
\newcommand{\E}{\mathbb{E}}
\newcommand{\eps}{\varepsilon}
\newcommand{\dd}{\,\mathrm d}
\newcommand{\sref}[2]{\hyperlink{source:#1:#2}{[S#2]}}
\newif\ifshowcatalogue
\showcataloguetrue
\newcommand{\cataloguescope}[1]{\ifshowcatalogue
 \paragraph{Statement recorded in the supplied source.}
 #1\fi}
\begin{document}
% UnsolvedMath ID 1097; source catalogue code GREEN-097

\setcounter{section}{1096}

\section{Asymptotic Enumeration of Nonattacking Queens}\label{1097}

{\small\sffamily UnsolvedMath ID 1097 · Combinatorics}

\subsection{Definitions and mathematical statement}

\paragraph{Shared notation.}Unless a section says otherwise, $\N=\{1,2,\ldots\}$,
$\Z$ is the ring of integers, $\Q,\R,\C$ are the rational, real, and complex
fields, $[N]=\{1,\ldots,\lfloor N\rfloor\}$, and $\log$ is the natural
logarithm. For real functions of a parameter tending to infinity,
$f\ll g$ and $f=O(g)$ mean $|f|\leq Cg$ eventually for some fixed $C$;
$f\gg g$ reverses this comparison; $f=o(g)$ means $f/g\to0$;
$f\sim g$ means $f/g\to1$; and $f\asymp g$ means both order bounds.
A subscript on $O$ or $\ll$ specifies allowed constant dependence.
The expression $n^{o(1)}$ in an upper bound means growth smaller than
$n^\epsilon$ for each fixed $\epsilon>0$. Local definitions govern any
exception, finite-field convention, or use of the same symbol for another
object. An asymptotic claim is not automatically an effective algorithm.



Represent a placement of $n$ queens by a permutation $\pi$ of $[n]$: the queens occupy $(i,\pi(i))$. Nonattack additionally requires the numbers $i-\pi(i)$ to be distinct and the numbers $i+\pi(i)$ to be distinct. Let $Q(n)$ count such permutations on the labelled board, without identifying rotations or reflections. The target is an asymptotic formula for this count. \sref{1097}{1} \sref{1097}{2}

\paragraph{Statement recorded in the supplied source.}
In how many ways (asymptotically) $Q(n)$ may $n$ non-attacking queens be placed on an $n \times n$ chessboard? \sref{1097}{1}

\paragraph{Scope and interpretation.}

The archive reports the leading logarithmic asymptotic as established. The residual question is an asymptotic enumeration beyond that precision; the known exponential-scale estimate is not being relabelled as open. \sref{1097}{1}

\paragraph{Residual task reported by the archive.}

The embedded review (2026-08-17) records: Obtain an asymptotic formula with multiplicative relative error, or otherwise determine the subexponential factor beyond Simkin's exponential rate. \sref{1097}{1}

\subsection{Short English statement}

How many ways can n mutually nonattacking queens be placed on a labelled n by n board, to asymptotically precise accuracy?

\subsection{Sources}

{\small

\begin{itemize}

\item[\hypertarget{source:1097:1}{S1}] ULAM.AI, \emph{UnsolvedMath}, supplied \texttt{problems.json}, record 1097 (GREEN-097), statement and background. The embedded literature-review date is 2026-08-17; that is the archive’s review date, not a fresh certification. Dataset landing page: \url{https://huggingface.co/datasets/ulamai/UnsolvedMath}.

\item[\hypertarget{source:1097:2}{S2}] Michael Simkin, The number of n-queens configurations, Advances in Mathematics 427 (2023), 109127. \url{https://doi.org/10.1016/j.aim.2023.109127}. Bibliographic information imported from the supplied record.

\item[\hypertarget{source:1097:3}{S3}] Parth Nobel, Akshay Agrawal, and Stephen Boyd, Computing tighter bounds on the n-queens constant via Newton's method, Optimization Letters 17 (2023), 1229-1240. \url{https://doi.org/10.1007/s11590-022-01933-2}. Bibliographic information imported from the supplied record.

\item[\hypertarget{source:1097:4}{S4}] Ben Green, 100 Open Problems, Problem 97, updates through December 2025;. \url{https://people.maths.ox.ac.uk/greenbj/papers/open-problems.pdf}. The author-hosted document was inspected on 27 September 2026; the archive’s local code is not a problem-number locator in this paper.

\end{itemize}

}

\label{end:1097}
\end{document}
